% Lösungen Einheit 1 von 4 – Dreiecksflächen (zusammenbau v0.1)
\einheitenkopf{Einheit 1 von 4 · Dreiecksflächen}

\erg{9}{a) halbe Grundseite mal Höhe; gebraucht werden die Grundseite und die senkrechte Höhe des Dreiecks}
\erg{10}{a) Höhe zu PQ – sie steht senkrecht auf der Grundseite \quad b) $|PQ| = 5$, $|QR| = 10$, $A = \tfrac{1}{2} \cdot 5 \cdot 10 = 25$ \quad c) $|PQ| = 6$, $M(3 | 2 | 1)$, $|MR| = 3$, $A = \tfrac{1}{2} \cdot 6 \cdot 3 = 9$ \quad d) $S_x(4 | 0 | 0)$, $S_y(0 | -6 | 0)$, Katheten 4 und 6: $A = \tfrac{1}{2} \cdot 4 \cdot 6 = 12$ \quad e) $P'(1 | 2 | 0)$, $Q'(4 | 2 | 0)$, $R'(1 | 5 | 0)$; rechtwinklig bei $P'$: $A = \tfrac{1}{2} \cdot 3 \cdot 3 = 4{,}5$ \quad f) $|OP| = 10$, die Höhe ist $z$ (R liegt über der Basismitte): $\tfrac{1}{2} \cdot 10 \cdot z = 40$, $z = 8$ \quad g) $T(-4 | 1 | 1)$; $RT$ ist parallel zu $PQ$ und dreimal so lang, gleiche Höhe: Trapez mit $\tfrac{1}{2} \cdot (a + 3a) \cdot h = 2ah$, Verhältnis $= \tfrac{1}{4}$ \quad h) S ist der Mittelpunkt von RT, denn T liegt von R aus doppelt so weit in Richtung RS. Die Dreiecke ORS und OST haben gleich lange Grundseiten RS und ST auf derselben Geraden und dieselbe Spitze O, also dieselbe Höhe – damit denselben Flächeninhalt.}
\erg{11}{a) Tom nimmt den Schenkel $|PR| = 5$ als Höhe. Die Höhe geht von $M(3 | 0 | 1)$ zu R: $|MR| = 4$, $A = \tfrac{1}{2} \cdot 6 \cdot 4 = 12$ \quad b) Die Höhe zerlegt das Dreieck in zwei rechtwinklige Dreiecke mit gleicher Hypotenuse (den Schenkeln) und gemeinsamer Kathete (der Höhe); nach Pythagoras sind auch die beiden Basisabschnitte gleich lang. \quad c) $A = 5 \cdot 4 - \tfrac{1}{2} \cdot 5 \cdot 1 - \tfrac{1}{2} \cdot 2 \cdot 4 - \tfrac{1}{2} \cdot 3 \cdot 3 = 9$ \quad d) $M(3 | 0 | 2)$, $|MR| = 5\,\mathrm{m}$, $A = \tfrac{1}{2} \cdot 6 \cdot 5 = 15\,\mathrm{m}^2$, Kosten $= 15 \cdot 12\,€ = 180\,€$}
